\chapter{1987 Nobel Prize in Economics}
\textbf{Laureates: }Robert Solow (USA)

\section*{Reason for Award}
For his contributions to the theory of economic growth

\section{What They Did}
Robert Solow developed the exogenous growth model showing capital accumulation
and technological progress drive long-run productivity growth. He identified the
Solow residual measuring unexplained technological progress.

His neoclassical growth model demonstrated conditional convergence, showing
poorer economies catch up with richer ones if save rates and technologies are
similar. He showed steady-state growth requires exogenous technological
advancement.

The Solow model became the standard teaching tool for growth economics. It
presents intuitive graphical analysis of capital accumulation, diminishing
returns, and transitional dynamics. Long-run growth is driven by technological
progress, not mere accumulation.

Solow's 1956 paper provided the first rigorous mathematical model of long-run
growth. It synthesized Keynesian macroeconomics with neoclassical production
theory. His 1957 paper introduced growth accounting framework decomposing output
growth into contributions from capital, labor, and total factor productivity.

\section{Impact on Economic Thinking}
Solow transformed growth economics from philosophical speculation into
quantitative analysis. He provided a framework where capital accumulation
contributes to growth but diminishing returns limit its effect. Technological
progress proved essential for sustained growth.

His identification of the residual attributed to technological change became the
standard growth accounting procedure. The conditional convergence hypothesis
explained patterns in cross-country incomes. Poorer nations tend to grow faster,
catching up with rich ones when saving rates and institutions are similar.

His model demonstrated the importance of savings rates and investment in capital
formation. It also emphasized ultimate dependency on technological advancement
for long-run growth. Growth accounting exercises worldwide decompose growth into
capital, labor, and total productivity components following Solow methodology.

The Solow model provided the first coherent framework for thinking about long-run
economic growth. It replaced earlier qualitative discussions with a rigorous
mathematical model. Its simplicity and transparency made it the standard
introductory tool for growth economics.

\section{Modern Perspectives Today}
Endogenous growth theories emerged challenging Solow's exogenous technological
progress assumption. They internalized innovation, knowledge spillovers, and
human capital accumulation into growth models. Yet the Solow model remains the
fundamental benchmark for growth analysis.

It is taught in introductory macroeconomics courses worldwide as the foundation
for understanding growth. Empirical studies testing convergence hypotheses use
the Solow framework as a starting point. Policy makers evaluating growth
strategies reference Solow conclusions about technological progress.

Research extending the Solow model incorporates environmental sustainability,
demographic transition, and climate change. These reflections demonstrate Solow's
enduring relevance. The fundamental insight that long-run growth depends on
technological progress rather than capital accumulation alone remains
uncontroversial.

The growth accounting methodology he developed remains the standard tool for
decomposing economic performance. It continues to inform debates about
productivity and living standards.
